\documentclass[12pt]{article}
\usepackage[utf8]{inputenc}
\usepackage[T1]{fontenc}
\usepackage[margin=1in]{geometry}
\usepackage{setspace}
\usepackage{xcolor}
\usepackage[colorlinks=true, urlcolor=blue!60!black, linkcolor=blue!60!black]{hyperref}
\setlength{\parindent}{2em}
\setlength{\parskip}{1.25em}
\renewcommand{\baselinestretch}{1.0}
\begin{document}
\begin{center}
{\LARGE \textit{United States v. Carolene Products Co.}, 304 U.S. 144 (1938)\par}
\vspace{0.8em}
{\large Case Brief --- Four-Page Edition\par}
\end{center}
\vspace{1em}

\section{Memory Jogger}

In upholding Congress's ban on interstate shipments of ``filled milk,'' the Court cemented deferential rational-basis review for ordinary economic regulation while planting, in Footnote Four, the blueprint for heightened scrutiny of laws infringing specific constitutional rights, distorting the political process, or targeting ``discrete and insular minorities'' (\href{https://supreme.justia.com/cases/federal/us/304/144/}{304 U.S. 144}).

\section{Detailed Case Facts}

The Filled Milk Act of 1923 barred interstate shipment of ``filled milk''---milk, cream, or skimmed milk blended with any fat or oil other than milk fat so as to imitate milk or cream---declaring it ``an adulterated article of food, injurious to the public health,'' whose sale ``constitutes a fraud upon the public.'' Congressional committees found that replacing butterfat (a vital vitamin source, especially for children) with cheaper vegetable oils produced a product that looked and tasted like condensed milk, invited deceptive marketing, and risked malnutrition. By 1938, roughly thirty-five states had banned filled milk.

Carolene Products shipped ``Milnut,'' a compound of condensed skimmed milk and coconut oil, in interstate commerce and was indicted under the Act. It argued that the Act exceeded the commerce power, invaded powers reserved to the states, violated Fifth Amendment due process by making Congress's declaration of harm conclusive, and arbitrarily spared oleomargarine.

\section{Procedural History}

Relying on its own 1934 decision holding the Act unconstitutional, the district court sustained the company's demurrer and dismissed the indictment. (The Seventh Circuit, in separate civil litigation, had upheld the Act.) The United States appealed directly to the Supreme Court, which reversed on April 25, 1938, holding that the demurrer should have been overruled.

\section{Judicial Votes}

Reversed, \textbf{6--1}, with seven Justices participating.
\begin{itemize}
\item \textbf{Opinion of the Court:} Justice Stone, joined in full by Chief Justice Hughes and Justices Brandeis and Roberts. Justice Black joined Parts First and Second only, so Part Third and Footnote Four drew four votes.
\item \textbf{Concurring in the result:} Justice Butler.
\item \textbf{Dissenting without opinion:} Justice McReynolds.
\item \textbf{Not participating:} Justices Cardozo and Reed.
\end{itemize}

\section{Holding}

The Filled Milk Act is valid on its face. Congress may exclude from interstate commerce articles it reasonably conceives to be injurious to public health, and regulatory legislation affecting ordinary commercial transactions is presumed constitutional unless the known facts preclude the assumption that it rests on some rational basis. The Fifth Amendment has no Equal Protection Clause, and a legislature may address one abuse without reaching all. Footnote Four, which suggests ``more exacting'' scrutiny for certain laws, is dicta that commanded only a four-Justice plurality.

\section{Analysis of Opinions}

\textbf{Part First.} The commerce power is ``complete in itself'' (\textit{Gibbons v. Ogden}) and includes excluding harmful articles from interstate commerce; exercising it in ways that resemble a state police power does not violate the Tenth Amendment.

\textbf{Part Second.} Relying on \textit{Hebe Co. v. Shaw} and the congressional record of nutritional harm and consumer confusion, Stone rejected the due process and equality challenges. Whether labeling would suffice was ``a matter for the legislative judgment, and not that of courts,'' and Congress could ban filled milk without banning oleomargarine.

\textbf{Part Third (four Justices).} The statutory declaration of harm is simply a legislative finding; facts supporting ordinary commercial legislation are presumed to exist, and where the legislative judgment is ``at least debatable,'' neither court nor jury may substitute its own view. Footnote Four identified three settings where the presumption may have ``narrower scope'': laws within a specific constitutional prohibition such as the Bill of Rights; laws restricting the political processes that ordinarily repeal undesirable legislation; and laws reflecting prejudice against discrete and insular religious, national, or racial minorities.

\textbf{Butler (concurring in the result).} The company must be allowed at trial to prove that Milnut was wholesome and not deceptive. \textbf{McReynolds} would have affirmed.

\section{Examples: Future Applications}

\begin{enumerate}
\item \textbf{Same-Side (Upheld): Licensing Interior Designers.} A state requires a degree and an exam for commercial interior designers, arguably to shield incumbents. It is ordinary economic regulation touching no specific right or discrete minority; because the safety rationale is ``at least debatable,'' it survives.
\item \textbf{Same-Side (Upheld): A Selective Butter-Substitute Ban.} Congress bans interstate shipment of a canola spread packaged to resemble butter while leaving margarine alone. Like filled milk, the legislature may address problems one step at a time.
\item \textbf{Opposite-Side (Heightened Scrutiny): Permits for Political Leaflets.} A city bans leaflets criticizing zoning laws without a mayor's permit. The ordinance falls within a specific First Amendment prohibition and clogs the political process, so the presumption narrows.
\item \textbf{Opposite-Side (Heightened Scrutiny): Excluding a Religious Minority.} A state bars businesses owned by members of a small, unpopular sect from public contracts. Footnote Four's protection of discrete and insular minorities calls for searching review.
\item \textbf{Fence-Sitter: Banning a Safe Milk Substitute.} A state bans clearly labeled, federally certified lab-grown dairy protein after hearings featuring only incumbent dairy farmers. Stone's test supplies a conceivable rational basis, but Part Third's caveat for changed facts and Butler's concurrence support a challenge.
\end{enumerate}

\section{Critique}

\textbf{Originalist and libertarian views.} Scholars such as Randy Barnett and Evan Bernick argue that the extreme presumption of constitutionality abdicates the judiciary's duty to guard against arbitrary, rent-seeking legislation; Geoffrey Miller's public-choice history portrays the Act as dairy-lobby protectionism against a wholesome competitor. Textualists add that nothing in the Constitution ranks personal liberties above economic rights, as Footnote Four's taxonomy does.

\textbf{Progressive and doctrinal views.} John Hart Ely built his theory of representation reinforcement on Footnote Four's second and third paragraphs. Bruce Ackerman responded that discrete and insular minorities may organize more effectively than diffuse groups such as consumers or the poor, and later scholars note that modern doctrine applies strict scrutiny symmetrically.

\section{Key Quotations}

\begin{quote}
``Even in the absence of such aids, the existence of facts supporting the legislative judgment is to be presumed, for regulatory legislation affecting ordinary commercial transactions is not to be pronounced unconstitutional unless, in the light of the facts made known or generally assumed, it is of such a character as to preclude the assumption that it rests upon some rational basis within the knowledge and experience of the legislators.''
\end{quote}
\noindent --- Justice Stone, Opinion of the Court, 304 U.S. at 152.

\begin{quote}
``There may be narrower scope for operation of the presumption of constitutionality when legislation appears on its face to be within a specific prohibition of the Constitution, such as those of the first ten amendments, which are deemed equally specific when held to be embraced within the Fourteenth.''
\end{quote}
\noindent --- Justice Stone, Footnote Four, 304 U.S. at 152 n.4.

\begin{quote}
``Nor need we enquire whether similar considerations enter into the review of statutes directed at particular religious, \dots{} or national, \dots{} or racial minorities, \dots{}: whether prejudice against discrete and insular minorities may be a special condition, which tends seriously to curtail the operation of those political processes ordinarily to be relied upon to protect minorities, and which may call for a correspondingly more searching judicial inquiry.''
\end{quote}
\noindent --- Justice Stone, Footnote Four, 304 U.S. at 152--53 n.4.

\end{document}
